% ==============================================================
\section{関連研究}\label{sec:related}
% ==============================================================

本研究は、時空間トランザクションデータから
密集アイテムセット領域を発見する問題を扱う。
この問題は2つの研究系譜の交点に位置する：
(1)~アイテムセットの列挙を目的とする\emph{パターンマイニング系譜}と、
(2)~多次元空間における密集ブロックの検出を目的とする\emph{密ブロック検出系譜}である。
以下では両系譜を整理し、本研究がなぜその交点に位置するかを明確にする。

% ==============================================================
\subsection{パターンマイニング系譜}
% ==============================================================

\paragraph{頻出アイテムセットマイニング（FIM）.}
Agrawal ら~\cite{agrawal1994apriori}の Apriori アルゴリズムは、
アイテムセットの単調性（部分集合は常に親集合以上のサポートを持つ）を利用して
候補空間を大幅に枝刈りする。
Han ら~\cite{han2000mining}の FP-Growth は
接頭辞木を用いてさらに高速化する。
これらは静的なトランザクションデータベース全体の
グローバルサポートを基準とするため、
局所的に密集したパターンを見逃すという
根本的な問題を抱える~\cite{aprioriwindow}。

\paragraph{局所密集区間検出手法.}
局所密集区間の検出を目的として、
連続する出現タイムスタンプ間の\emph{間隔}に基づく
複数の手法が提案されている。
LPFIM~\cite{mahanta2005lpfim}は
連続する2出現のギャップが閾値 minthd1 を超えた時点で区間を分割し、
局所サポート率が閾値以上かつ区間長が minthd2 以上の区間を出力する。
LPPM~\cite{fournierviger2021lppm}は
余剰累積（spillover）を導入してノイズに頑健な区間検出を実現する。

これらの手法は
「連続する2出現の時刻差（出現間隔）」を評価することで密集を推定する。
しかし、高橋ら~\cite{aprioriwindow}が指摘するように、
ギャップパラメータがパターン同定精度と区間検出精度の双方を支配するため、
両者を同時に高く保つパラメータ設定が一般に存在しない
という構造的限界を持つ。
さらに、これらの手法は\emph{本質的に多次元への拡張が不可能}である。
$d$ 次元グリッド上の出現点集合には線形順序が存在せず、
「連続する2出現の間隔」が定義できないためである。

Apriori-Window~\cite{aprioriwindow}は
この構造的限界を解消する厳密解法である。
出現間隔を介さず、スライディング窓 $[l, l+W]$ 内の出現数を\emph{直接}計算することで
局所頻度を評価するため、ギャップパラメータを必要としない。
反単調性に基づく Apriori 枝刈りと $t$ 軸ストライドにより、
密集区間を網羅的かつ厳密に列挙する。
比較実験では LPFIM・LPPM を大幅に上回る性能を示した~\cite{aprioriwindow}。

\emph{多次元化が可能な理由}：
「矩形窓内の出現数」という定義は $d$ 次元でも自然に拡張できる
（窓を $d$ 次元直方体 $\prod_k [s_k, s_k+W_k]$ に置き換えるだけ）。
出現間隔ベースの手法が線形順序に依存するのと対照的に、
Apriori-Window の設計は線形順序を必要としない。
本研究はこの観察に基づき、
Apriori-Window を1次元から任意の $d$ 次元へ拡張する。

\paragraph{コロケーションマイニング.}
コロケーションマイニング~\cite{shekhar2001colocation}は、
複数種のフィーチャーが地理的近傍に共存するパターンを
参加指数（PI）で評価する。
Qian ら~\cite{qian2009stcolocation}はこれを
重み付きスライディングウィンドウに拡張した（STCP-Miner）。
これらはトランザクション単位の共起ではなく
空間的近傍に基づく共起を扱い、
密集領域の抽出は対象としない。
ただし、パターン識別精度（F1）の観点では比較が可能である（\S\ref{sec:ex1_1}）。

% ==============================================================
\subsection{領域抽出系譜}
% ==============================================================

% --- 密ブロック検出 ---
\paragraph{密ブロック検出系手法.}
M-Zoom~\cite{shin2018mzoom}をはじめとする
密部分テンソル手法群（テンソル:~\cite{jiang2015crossspot,shin2018mbiz,shin2017dcube,shin2017densealert,shin2023catchcore}、
グラフ: SpecGreedy~\cite{feng2020specgreedy}）は、
多次元データにおける密ブロックを検出する。
各次元で値の集合 $B_k \subseteq D_k$ を選択し、
直積ブロック $\prod_k B_k$ の質量/体積比を最大化する
グリーディ探索を行う。

これらの手法には2つの構造的限界がある。

\emph{第一に、出力が直積ブロックに限定される。}
不規則な形状の密集領域に対しては、
過剰な領域を含むか密集セルを取りこぼすかが不可避であり、
GT パターンを与えても Jaccard が低下する
（\S\ref{sec:ex1_1}で実証）。

\emph{第二に、パターン列挙が不正確である。}
アイテムを追加次元として含めれば列挙は試みられるが、
直積制約によりブロック内のアイテム集合 $C_B$ の全部分集合が
同一領域で密集と仮定され、偽陽性が生じる（\S\ref{sec:ex1_1}で実証）。

% --- 空間統計手法 ---
\paragraph{空間統計手法.}
LISA~\cite{anselin1995lisa}は局所 Moran's $I$ で空間ホットスポットを同定し、
Getis-Ord $G_i^*$~\cite{ord1995getis}は
局所集中度を $z$ スコアで評価する。
いずれも\emph{空間次元のみ}が対象であり、
時間スライスごとの個別実行では時間的連続性が失われる。

% --- 空間・時空間スキャン統計 ---
\paragraph{空間・時空間スキャン統計.}
Kulldorff~\cite{kulldorff1997spatial,kulldorff2001prospective}の
空間/時空間スキャン統計は、
円形・円柱形の走査窓で有意なクラスタを検出する。
FlexScan~\cite{tango2005flexscan}は
不整形窓への拡張を行った。

上記の統計手法はいずれも変量ごとの異常集中を検出する設計であり、
トランザクション内の共起パターン（アイテムセット）を列挙する機能を持たない。
全パターンの列挙には $O(2^{|\mathcal{I}|})$ 回の独立実行が必要であり、
Apriori 枝刈りとの統合は設計されていない。

% --- 時空間クラスタリング ---
\paragraph{時空間密度ベースクラスタリング（ST-DBSCAN 等）.}
ST-DBSCAN~\cite{birant2007stdbscan}は
DBSCAN~\cite{ester1996dbscan}を時空間データに拡張し、
空間的距離と時間的距離にそれぞれ独立な閾値
$\varepsilon_1, \varepsilon_2$ を設ける。
任意形状のクラスタを出力できるが、
第一に、点近傍のカウントにより密度を定義するため、
出現が空間的に均一な場合にはデータ全体が単一クラスタとなり、
局所的密集を識別できない。
第二に、点近傍密度は反単調性を持たないため、
全パターンの独立処理 $O(2^{|\mathcal{I}|})$ が不可避である。

上記の領域抽出系手法はいずれもアイテムセットの概念を持たないが、
GT パターンとその出現密度を与えれば
領域抽出器として評価が可能である（\S\ref{sec:ex1_1}）。

% ==============================================================
\subsection{本研究の位置づけ}
% ==============================================================

以上の分析により、既存手法はいずれも
パターン列挙と多次元密集領域抽出を同時に達成できないことが示された。
本研究は窓内出現数に基づく密集定義の反単調性を利用して
Apriori 枝刈りを自然に統合し、
両者を単一のフレームワークで実現する。
表~\ref{tab:related}にこの位置づけを示す。

\paragraph{より強い比較手法が存在しない理由.}
表~\ref{tab:related}の空白は手法選択の問題ではなく、
2つの列間の\emph{構造的}非両立性を反映している。
パターンマイニング系手法は反単調な密集領域を持たないため
時空間ウィンドウを出力できず、
領域抽出系手法はアイテムセットの概念を持たないため
全$2^{|\mathcal{I}|}$部分集合を独立に処理する必要があり
Apriori枝刈りと統合できない。
同一カテゴリ内の高速化手法
（例：より高速なコロケーションマイナー、より高速な密部分テンソル探索）は
同じ構造的制約を共有しており、
\S\ref{sec:ex1}で示す根本的なミスマッチは変わらない。
実験のbaselineは各手法カテゴリを代表する手法を網羅的に選択しており、
本研究が埋めるギャップを包括的に特徴付けている。

\begin{table}[t]
\centering
\caption{関連手法との比較。
「アイテムセット」はパターン列挙能力、
「密集領域」は時空間的に局所化された密集領域の出力能力を指す。}
\label{tab:related}
\small
\begin{tabular}{lccc}
\toprule
手法 & \shortstack{アイテム\\セット} & \shortstack{時空間\\密集領域} & \shortstack{領域\\形状} \\
\midrule
Apriori~\cite{agrawal1994apriori}
                                          & $\checkmark$ & $\times$     & ---          \\
LPFIM~\cite{mahanta2005lpfim} / LPPM~\cite{fournierviger2021lppm}
                                          & $\checkmark$ & 1Dのみ      & 区間     \\
Apriori-Window~\cite{aprioriwindow}       & $\checkmark$ & 1Dのみ      & 区間     \\
コロケーション~\cite{shekhar2001colocation}
                                          & $\checkmark$ & $\times$     & ---          \\
STCP-Miner~\cite{qian2009stcolocation}   & $\checkmark$ & $\times$     & ---          \\
\midrule
M-Zoom~\cite{shin2018mzoom} 他           & $\triangle$$^*$     & $\checkmark$ & 直積のみ \\
SaTScan~\cite{kulldorff1997spatial}       & $\times$     & $\checkmark$ & 円/円柱 \\
FlexScan~\cite{tango2005flexscan}        & $\times$     & $\checkmark$ & 不整形    \\
LISA~\cite{anselin1995lisa}
  / $G_i^*$~\cite{ord1995getis}          & $\times$     & 空間のみ     & 任意         \\
ST-DBSCAN~\cite{birant2007stdbscan} 等   & $\times$   & $\checkmark$ & 任意    \\
\midrule
\textbf{本手法}                              & $\checkmark$ & $\checkmark$ & 任意    \\
\bottomrule
\multicolumn{4}{l}{\scriptsize $^*$アイテムを次元に含めればパターン出力は可能だが、} \\
\multicolumn{4}{l}{\scriptsize \hspace{1em}直積制約により偽陽性が生じ領域分離も不可能。}
\end{tabular}
\end{table}
